\documentclass[12pt, ukrainian]{article}
\usepackage[a4paper]{geometry}
\setlength\parindent{0mm}
\geometry{top=1.0cm,bottom=1.0cm,left=1.3cm,right=1.3cm}
\pagestyle{empty}
\usepackage[T2A]{fontenc}
\usepackage[utf8]{inputenc}
\usepackage[ukrainian]{babel}
\usepackage{graphicx}
\usepackage{caption}
\usepackage{caption}
\DeclareCaptionFormat{nocaption}{}
\captionsetup{format=nocaption,aboveskip=0pt,belowskip=0pt}
\usepackage[Export]{adjustbox}
\adjustboxset{max size={0.9\linewidth}{0.9\paperheight}}
\usepackage{verbatim, listings, clrscode3e, fancyvrb}
\def\code#1{\Verb!#1!}
\begin{document}
%begin
\begin {large}

\textbf{\_\_.11.2025 Контрольна робота, варіант  \No { 144 }}\par


{

\begin{center}
\textbf{Завдання 1}
\end{center}
Нехай int var=13; Один потік збільшив змінну var на 3, а інший збільшив її в три рази.

гонка даних (data race), стан гонки (race condition)

Виберіть усі правильні твердження (якщо такі є).\par
\vspace{2mm}
(\,1\,) Тут гарантовано є гонка даних.%bad

(\,2\,) Тут може бути тільки гонка даних.%bad

(\,3\,) Тут гарантовано нема ані гонки даних, ані стану гонки.%bad

(\,4\,) Тут може не бути гонки даних.%good



\begin{center}
\textbf{Завдання 2}
\end{center}
Виберіть усі правильні твердження (якщо такі є).\par
\vspace{2mm}

(\,1\,) У стані гонки програма відчайдушно намагається захопити ресурси процесора.%bad

(\,2\,) Об’єкти класів promise та future можна використовувати тільки в багатопотокових програмах.%bad

(\,3\,) Значення, що є результатом виконання функції, запущеної через async, зберігається безпосередньо в об’єкті ф’ючерса (future).%bad

(\,4\,) Чи запустить функція  async виконання в новому потоці можна задавати політикою.%good



\begin{center}
\textbf{Завдання 3}
\end{center}
Виберіть усі правильні твердження (якщо такі є).\par
\vspace{2mm}

(\,1\,) За використання атомарних змінних структура даних є структурою даних без блокувань.%bad

(\,2\,) Структура даних є структурою даних без блокувань, якщо її реалізація не використовує м’ютекси.%bad

(\,3\,) Структура даних без блокувань гарантує, що жоден потік не буде очікувати доступу до виконання операцій над цією структурою даних.%bad

(\,4\,) За використання структурою даних рівно одного м’ютексу, взаємоблокування на її операціях неможливі.%good



\begin{center}
\textbf{Завдання 4}
\end{center}
Виберіть усі правильні твердження (якщо такі є).\par
\vspace{2mm}

(\,1\,) Потік може мати доступ до автоматичних змінних іншого потоку, але це небезпечно.%good

(\,2\,) Динамічні змінні не можуть розділятися між потоками.%bad

(\,3\,) Потік може мати доступ до автоматичних змінних іншого потоку і це безпечно.%bad

(\,4\,) Автоматичні змінні не можуть розділятися між потоками.%bad



\begin{center}
\textbf{Завдання 5}
\end{center}
У цьому фрагменті коду
\begin{verbatim}
template <class T> void f(T a){}

void g(){
    int i=5;
    f(unique_ptr<int>(new int(i)))+(unique_ptr<int>(new int(++i)));
\end{verbatim}
(\,1\,) є гонка даних%bad

(\,2\,) можливий виток ресурсів%good

(\,3\,) є невизначена поведінка%good

(\,4\,) є стан гонки%bad


}
\end {large}


\end{document}

